% ===========================================================================
% SOURCES: docs/claude_project/03_HARDWARE_ARCHITECTURE.md
%          docs/controller_design.md (Greek - the architecture study)
%          docs/pointwise_datapath_explained.md
%          docs/latency_map.md
%          docs/design_decisions.md: DD-013..DD-018, DD-021, DD-029
%          hardware/rtl/**
%
% This chapter describes WHAT WAS BUILT and why. It does NOT contain the
% timing-closure narrative (\cref{ch:timing}) or the verification method
% (\cref{ch:verification}) - both were split out so this chapter stays
% readable. Resist pulling them back in.
% ===========================================================================
\chapter{Accelerator Architecture}
\label{ch:architecture}

\section{Proposed Accelerator Architecture}
\label{sec:arch-overview}
% The block diagram, and the one-paragraph summary of the dataflow. A reader
% should be able to stop after this section and still know what the thing is.
% FIGURE: a top-level block diagram. Does not exist yet - thesis/figures/ has
% only pe_array.svg. This is the single most valuable figure in the thesis.

\section{Design Philosophy: Engines Do Only Arithmetic}
\label{sec:engines-philosophy}
% Set on the first RTL day and never broken: the engines take operands and
% produce results; addressing, windowing and sequencing live in the feeders and
% the controller. Two consequences worth stating, because they pay off later:
% every engine is unit-testable against a golden vector with no controller
% present (\cref{ch:verification}), and the integration testbenches that build
% windows by hand became the specification of the controller.

\section{Parallelizing the Compute Engines}
\label{sec:parallelism}

\subsection{Pointwise Convolution: A \texorpdfstring{$32 \times 16$}{32x16} MAC Array}
\label{sec:pw-array}
% \DD{013}. Tm=32 output lanes x Tn=16 input channels = 512 MAC/cycle, chosen
% against the tiling space of \cref{sec:tiling} and the \thPwMacPct\% MAC share
% of \cref{sec:mac-dist}. Output-stationary; accumulation over ceil(IC/Tn)
% cycles. Broadcast activations, per-lane weights.
\begin{figure}[t]
  \centering
  % figures/pe_array.pdf is produced from the SVG by `make figures', which
  % needs Inkscape. The guard keeps the manuscript compiling on a machine that
  % does not have it yet - a missing figure should not stop you writing.
  \IfFileExists{figures/pe_array.pdf}{%
    \includegraphics[width=\linewidth]{pe_array.pdf}%
  }{%
    \fbox{\parbox{0.8\linewidth}{\centering\vspace{2em}
      \texttt{figures/pe\_array.pdf} not built --- run \texttt{make figures}
    \vspace{2em}}}%
  }
  \caption{The pointwise MAC array. The $T_m = 32$ lanes share one broadcast of
           $T_n = 16$ activations while each lane holds its own weights, giving
           $P = T_m \cdot T_n = 512$ MAC per cycle.}
  \label{fig:pe-array}
\end{figure}

\subsection{Depthwise Convolution: A 16-Channel Engine}
\label{sec:dw-engine}
% \DD{014}, Tc=16. Separate from the pointwise array because the dataflow is
% different, not because of the MAC count. The line buffer reads each input
% pixel ONCE rather than nine times - that is the result worth reporting here.

\subsection{Stem Convolution: An 8-Channel Engine}
\label{sec:stem-engine}
% Runs once per inference, on 3 input channels. Sized accordingly; justify the
% asymmetry with \thStemMacPct\%.

\subsection{Requantization Lanes}
\label{sec:requant-lanes}
% The shared out_stage. \DD{016}: sharing it took the predicted DSP count from
% ~634 to 442. Implements \cref{sec:requant-theory} exactly.

\subsection{Datapath Widths}
\label{sec:widths}
% Every width, and for each one whether it is a BOUND or a MEASUREMENT. After
% \DD{019} this table is one of the more defensible things in the thesis -
% make the distinction a column, not a footnote.

\section{Memory Organization}
\label{sec:memory}

\subsection{On-Chip Activations and Off-Chip Weights}
\label{sec:mem-split}
% Activations never leave the chip; weights stream from DDR. Justify from
% \cref{sec:roofline-analysis}: \thWeightMB{} MB of weights against activations
% that fit in URAM.

\subsection{Activation Pool and Static Allocation}
\label{sec:act-pool}
% One pool, statically allocated by the compiler. No allocator at run time, so
% no fragmentation and no dynamic failure mode.

\subsection{Weight, Parameter and Image Buffers}
\label{sec:buffers}
% Sizes, banking, and the address-width headroom noted in \cref{sec:risks}:
% features.18.0 uses 800 of 1024 entries, and the classifier is small only
% because this model has \thClasses{} classes.

\section{Control Unit and Instruction Set}
\label{sec:control}

\subsection{Program Sequencer}
\label{sec:sequencer}
% Sequential issue, one op at a time, \thInstructions{} instructions. State the
% ceiling plainly: 64 of 64 used, zero spare, and it fails at build time rather
% than corrupting execution.

\subsection{Instruction Format and Program Compiler}
\label{sec:isa-format}
% scripts/gen_program.py compiles manifest.json to the program. Full encoding
% in \cref{ch:isa}; here, only the shape and the six opcodes.

\subsection{Feeders and Handshaking}
\label{sec:feeders}
% Five feeders, one per dataflow. The addr_gen contract - shape latched at
% start so the caller may prefetch - and the C-3 finding: no feeder honoured
% it, masked only because top_seq does not refetch until !op_busy. Worth
% telling, because implementing the advertised prefetch would have broken all
% five silently, as wrong addresses rather than as an error.

\subsection{Resource Sharing}
\label{sec:sharing}
% \DD{016}, \DD{017}. out_stage shared across engines; res_feeder needs no
% multiplier; gap_feeder needs no DSP; the classifier needs no feeder at all.
% \DD{017} - two busy bugs that appeared only when sharing became real - is
% evidence for \cref{sec:three-levels} and should be cross-referenced there.

\section{Latency as a Contract}
\label{sec:latency-contract}
% SOURCE: docs/latency_map.md. Pipeline depths are restated in five feeders, so
% a kernel cannot change in isolation. This is the structural fact that made
% the timing work in \cref{ch:timing} expensive, and A-2 (dw_array exporting
% valid/valid_out so latency is a wire rather than a constant) is the fix
% pattern. Put the latency table here and refer to it from \cref{ch:timing}.
